\documentclass[12pt,oneside]{article}

%%%%%%%%%%%%%%%%%%%%%%%%%%%%
%%   Zusaetzliche Pakete  %%
%%%%%%%%%%%%%%%%%%%%%%%%%%%%
\usepackage{enumerate}  
\usepackage{fancyhdr}
\usepackage{a4wide}
\usepackage{graphicx}
\usepackage{palatino}
\usepackage{multirow}
\usepackage{booktabs}
\usepackage{titlesec}
\usepackage{acronym}% http://ctan.org/pkg/acronym
\usepackage{enumitem}% http://ctan.org/pkg/enumitem

%folgende Zeile auskommentieren für englische Arbeiten
%\usepackage[ngerman]{babel}
%folgende Zeile auskommentieren für deutsche Arbeiten
\usepackage[ngerman, english]{babel}

\usepackage[T1]{fontenc}
\usepackage[utf8]{inputenc}
\usepackage[bookmarks]{hyperref}
\usepackage[justification=centering]{caption}
\usepackage[style=apa,natbib=true,backend=biber,maxbibnames=20]{biblatex}
\usepackage{csquotes}
\bibliography{literature}

\setlength{\parindent}{0em} 
\setlist[itemize]{noitemsep, topsep=0pt}
\setlist[enumerate]{noitemsep, topsep=0pt}


%%%%%%%%%%%%%%%%%%%%%%%%%%%%%%
%% Definition der Kopfzeile %%
%%%%%%%%%%%%%%%%%%%%%%%%%%%%%%

\pagestyle{fancy}
\fancyhf{}
\cfoot{\thepage}
\setlength{\headheight}{16pt}

%%%%%%%%%%%%%%%%%%%%%%%%%%%%%%%%%%%%%%%%%%%%%%%%%%%%%
%%  Definition des Deckblattes und der Titelseite  %%
%%%%%%%%%%%%%%%%%%%%%%%%%%%%%%%%%%%%%%%%%%%%%%%%%%%%%

\newcommand{\JMUTitle}[9]{

  \thispagestyle{empty}
  \vspace*{\stretch{1}}
  {\parindent0cm
  \rule{\linewidth}{.7ex}}
  \begin{flushright}
    \vspace*{\stretch{1}}
    \sffamily\bfseries\Huge
    #1\\
    \vspace*{\stretch{1}}
    \sffamily\bfseries\large
    #2\\
    \vspace*{\stretch{1}}
    \sffamily\bfseries\small
    #3
  \end{flushright}
  \rule{\linewidth}{.7ex}

  \vspace*{\stretch{1}}
  \begin{center}
    \includegraphics[width=2in]{siegel} \\
    \vspace*{\stretch{1}}
    \Large Exposé  \\

    \vspace*{\stretch{2}}
   \large Lehrstuhl f\"{u}r Wirtschaftsinformatik\\
    \large und Systementwicklung\\
    \large Universität Würzburg\\
    \vspace*{\stretch{1}}
    \large Betreuer:  #8 \\[1mm]
    
    \vspace*{\stretch{1}}
    \large W\"urzburg, den #7 \\
        \vspace*{\stretch{0.25}}

    % Bearbeitungszeit: 14.03.2025 - 09.05.2025 % Die Bearbeitungszeit der Seminar-/ Abschlussarbeit ist hier einzutragen.

  \end{center}
}

\titlespacing*{\section}
{0pt}{3.5ex plus 1ex minus .2ex}{.2ex plus .2ex}
\titlespacing*{\subsection}
{0pt}{1.5ex plus 1ex minus .2ex}{.2ex plus .2ex}
\titlespacing*{\subsubsection}
{0pt}{1.5ex plus 1ex minus .2ex}{.2ex plus .2ex}




%%%%%%%%%%%%%%%%%%%%%%%%%%%%
%%  Beginn des Dokuments  %%
%%%%%%%%%%%%%%%%%%%%%%%%%%%%

\begin{document}

  \JMUTitle
      {Data Objects in ML Contexts: A Systematic Literature Review in IS Research}        % Titel der Arbeit
      {Linus Schärmann}                        % Vor- und Nachname des Autors
      {2910412}
      
      {Wirtschaftswissenschaftlichen Fakultät}  % Name der Fakultaet
      {W"urzburg 2025}                          % Ort und Jahr der Erstellung
      {08.11.2025}                              % Tag der Abgabe
      {Tim Thorwart-Gumpert}               % Name des Erstgutachters
      {}                          % Name des Zweitgutachters

  \clearpage

\lhead{}
\pagenumbering{Roman} 
    \setcounter{page}{1}

\tableofcontents
\clearpage

%%%%%%%%%%%%%%%%%%%%%%%%%%%%
%%  Kurzzusammenfassung   %%
%%%%%%%%%%%%%%%%%%%%%%%%%%%%
% \newpage
% \lhead{Abstract}
% \section*{Abstract}
% \addcontentsline{toc}{section}{Abstract}

%\newpage
%\lhead{List of Figures} % Bei englischsprachiger Arbeit anzupassen auf: List of Figures
%\addcontentsline{toc}{section}{List of Figures} % Bei englischsprachiger Arbeit anzupassen auf: List of Figures
%\listoffigures

% \newpage
% \lhead{List of Tables} % Bei englischsprachiger Arbeit anzupassen auf: List of Tables
% \addcontentsline{toc}{section}{List of Tables} % Bei englischsprachiger Arbeit anzupassen auf: List of Tables
% \listoftables
% \newpage

\setlength{\parskip}{0.5em} 


%%%%%%%%%%%%%%%%%%%%%%%%%%%%%%%%%%
%%  Definition der Abkürzungen  %%
%%%%%%%%%%%%%%%%%%%%%%%%%%%%%%%%%%
% \lhead{List of Abbreviations} % Bei englischsprachiger Arbeit anzupassen auf: List of Abbreviations
% \section*{List of Abbreviations} % Bei englischsprachiger Arbeit anzupassen auf: List of Abbreviations
% \addcontentsline{toc}{section}{List of Abbreviations} % Bei englischsprachiger Arbeit anzupassen auf: List of Abbreviations

% \begin{acronym}

% \end{acronym}

%%%%%%%%%%%%%%%%%%%%%%%%%%%%
%%  Einstellungen  %%
%%%%%%%%%%%%%%%%%%%%%%%%%%%%
\clearpage
\pagenumbering{arabic}  
    \setcounter{page}{1}
\lhead{\nouppercase{\leftmark}}

%%%%%%%%%%%%%%%%%%%%%%%%%%%%
%%  Hauptteil  %%
%%%%%%%%%%%%%%%%%%%%%%%%%%%%

\section{Research Topic} \label{researchtopic}
Short introduction, explanation of the topic and its relevance. 

\section{Objectives and Research Interests} \label{objectives}
Explain what you want to find out by the end of your research project—and why!

\section{State of Research and Theoretical Basis} \label{stateresearch}
A brief (not too detailed) overview of the existing literature and — depending on the field of research — a theoretical basis for your research.

\section{Concept} \label{concept}
The concept is the core of your exposé: Here you present your question, your hypothesis (assumed result), and, if applicable, your methodology.

\section{Preliminary Outline}
Even if your outline isn't quite ready yet, this is a good place to structure your initial ideas.

\section{Personal Motivation or Preparatory Work} \label{motivation}
Optional – If you have already worked on this topic or have specific motivations for your research, you can describe them here.

%%%%%%%%%%%%%%%%%%%%%%%%%%%%
%% Literaturverzeichnis wird 
%% automatisch eingefügt
%%%%%%%%%%%%%%%%%%%%%%%%%%%%
\clearpage
\lhead{}
\printbibliography
\addcontentsline{toc}{section}{\bibname}


%%%%%%%%%%%%%%%%%%%%%%%%%%%%
%% Anhang (optional) 
%%%%%%%%%%%%%%%%%%%%%%%%%%%%
%\clearpage
%\appendix
%\section{Appendix A} % Bei englischsprachiger Arbeit anzupassen auf: Appendix A

%%%%%%%%%%%%%%%%%%%%%%%%%%%%
%% Eidesstattliche Erklärung
%% muss angepasst werden 
%% in Erklaerung.tex
%%%%%%%%%%%%%%%%%%%%%%%%%%%%
\newpage
\begin{otherlanguage}{ngerman}
\thispagestyle{empty}
\section*{Eidesstattliche Erklärung}
\thispagestyle{empty}
Hiermit versichere ich, die vorliegende Arbeit selbstständig verfasst und keine anderen als die angegebenen Quellen und Hilfsmittel benutzt sowie die Zitate deutlich kenntlich gemacht zu haben.
\newline
Ich erkläre weiterhin, dass die vorliegende Arbeit in gleicher oder ähnlicher Form noch nicht im Rahmen eines
anderen Prüfungsverfahrens eingereicht wurde.
\vspace{4\baselineskip}\\
Würzburg, den \today \hfill Linus Schärmann 
\vspace{4\baselineskip}\\
\end{otherlanguage}

\end{document}